\documentclass[11pt]{article}
\usepackage[utf8]{inputenc}
\usepackage[T1]{fontenc}
\usepackage[margin=2.5cm]{geometry}
\usepackage{graphicx}
\usepackage{amsmath}
\usepackage{booktabs}
\usepackage{array}
\usepackage{titlesec}
\usepackage{enumitem}
\usepackage{xcolor}
\usepackage{hyperref}

\titleformat{\section}{\normalfont\Large\bfseries}{\thesection}{1em}{}
\titleformat{\subsection}{\normalfont\large\bfseries}{\thesubsection}{1em}{}

\title{Bitácora de Diseño e Implementación\\
\large Proyecto Individual --- Lógica Combinacional: Puerta con Contraseña\\
\large CE 1107 --- Fundamentos de Arquitectura de Computadores}
\author{[Nombre del estudiante]}
\date{}

\begin{document}
\maketitle
\vspace{-1.5cm}

\begin{center}
\begin{tabular}{ll}
\textbf{Curso:} & CE 1107 --- Fundamentos de Arquitectura de Computadores \\
\textbf{Profesor:} & Luis Chavarría Zamora \\
\textbf{Institución:} & Instituto Tecnológico de Costa Rica \\
\textbf{Repositorio:} & [enlace al repositorio Git] \\
\end{tabular}
\end{center}

\vspace{0.5cm}
\hrule
\vspace{0.5cm}

\section*{Jueves 3 de septiembre de 2026}

\subsection*{9:00 a.m. -- 11:00 a.m. --- Investigación y decisiones de diseño}

Se dedicó este bloque a investigación preliminar y a la toma de decisiones fundamentales
de diseño antes de iniciar cualquier esquemático.

\textbf{Decisión de tecnología:} se optó por trabajar con compuertas \textbf{CMOS}
(familia 4000), en lugar de TTL, dado que CMOS ofrece mayor accesibilidad y comodidad
de uso en aplicaciones de baja potencia, además de mayor inmunidad al ruido y menor
consumo estático en comparación con TTL. Esta decisión condiciona todo el diseño
posterior, ya que fija los niveles de tensión de operación (0V--5V) y los
circuitos integrados específicos a utilizar.

Se investigaron los siguientes temas para sustentar el diseño:
\begin{itemize}[nosep]
    \item Funcionamiento interno y patrones de excitación de un display de
    7 segmentos (configuración cátodo común vs. ánodo común y su efecto sobre
    la lógica de manejo de cada segmento).
    \item Comportamiento de compuertas lógicas básicas en tecnología CMOS
    (inversor, AND, OR, XNOR) y sus equivalentes disponibles en la serie 4000
    (CD4069, CD4081, CD4073).
    \item Funcionamiento del flip-flop tipo D, con énfasis en su configuración
    como \emph{toggle} (realimentando $Q'$ hacia $D$) para lograr una salida que
    cambia de estado en cada pulso de reloj, sin necesidad de lógica de detección
    de flanco adicional.
\end{itemize}

Esta investigación fue la base conceptual utilizada en el bloque de trabajo de
la tarde para plantear los esquemáticos.

\subsection*{2:00 p.m. -- 4:00 p.m. --- Diseño de los esquemáticos por módulo}

Con base en la investigación de la mañana, se definió la arquitectura del sistema
en tres módulos independientes, cada uno con su propia lógica booleana:

\begin{enumerate}
    \item \textbf{Ingreso y verificación de contraseña.} Recibe la contraseña de
    3 dígitos en base 5 (código elegido: \texttt{1-4-3}), representada mediante
    9 líneas de entrada (3 bits por dígito), y produce una señal \texttt{Correcto}
    que indica si la combinación ingresada coincide con la contraseña definida.

    \item \textbf{Cambio de estado de la puerta.} Un flip-flop tipo D configurado
    como \emph{toggle}, que cambia el estado de la puerta (abierta/cerrada) cada
    vez que se confirma una contraseña correcta. El estado se indica mediante dos
    LEDs (verde = abierta, rojo = cerrada).

    \item \textbf{Anuncio en el display.} Un decodificador simplificado que
    muestra en un display de 7 segmentos si el sistema está \emph{leyendo}
    (símbolo ``L'') o si la contraseña ingresada fue \emph{inválida}
    (símbolo ``E'').
\end{enumerate}

\subsubsection*{Ecuaciones del módulo 1 --- Comparador de contraseña}

Para la contraseña \texttt{1-4-3} en base 5, cada dígito se compara contra su
valor binario fijo mediante una simplificación: dado que el valor de referencia
de cada bit ($B_i$) es constante, la compuerta XNOR de comparación se reduce a
una línea directa (si $B_i=1$) o a un inversor (si $B_i=0$).

\begin{align*}
\text{Dígito 1 (valor 1, código 001)}: \quad Eq_1 &= A_2' \cdot A_1' \cdot A_0 \\
\text{Dígito 2 (valor 4, código 100)}: \quad Eq_2 &= A_2 \cdot A_1' \cdot A_0' \\
\text{Dígito 3 (valor 3, código 011)}: \quad Eq_3 &= A_2' \cdot A_1 \cdot A_0 \\[4pt]
\text{Señal de contraseña correcta}: \quad Correcto &= Eq_1 \cdot Eq_2 \cdot Eq_3
\end{align*}

\subsubsection*{Ecuaciones del módulo 2 --- Disparo del toggle}

\begin{align*}
Trigger &= P \cdot Correcto
\end{align*}

donde $P$ es la señal del botón de confirmar. Esta señal alimenta la entrada de
reloj (CLK) de un flip-flop D en modo \emph{toggle} ($D = Q'$), de forma que el
estado de la puerta conmuta con cada confirmación correcta, sin requerir lógica
adicional de detección de flanco.

\subsubsection*{Ecuaciones del módulo 3 --- Anuncio en display}

\begin{align*}
S_0 &= P \cdot Correcto' \\
a = g &= S_0 \\
b = c &= 0 \text{ (GND fijo)} \\
d = e = f &= 1 \text{ (Vcc fijo)}
\end{align*}

\textbf{Decisión de diseño documentada:} originalmente el enunciado contemplaba
cuatro estados de anuncio (leyendo, inválida, abriendo, cerrando), lo que
requeriría 2 bits de estado. Sin embargo, dado que el estado de apertura/cierre
de la puerta ya se comunica mediante los LEDs del módulo 2, se decidió no
duplicar esa información en el display, reduciendo el sistema a solo 2 estados
(leyendo/inválida) representables con 1 bit ($S_0$). Esto permite que los
segmentos $b$, $c$, $d$, $e$ y $f$ queden cableados de forma fija a Vcc/GND, sin
necesidad de compuertas adicionales, simplificando considerablemente el circuito
sin pérdida de funcionalidad requerida para este avance.

\vspace{0.3cm}
\hrule
\vspace{0.3cm}

\section*{Sábado 5 de septiembre de 2026}

\subsection*{9:00 a.m. -- 1:00 p.m. --- Implementación en TinkerCAD}

Se implementó el diseño de los tres módulos en el simulador TinkerCAD, elegido
por su simplicidad de uso y por ofrecer una experiencia de armado relativamente
cercana a la de un circuito físico real en protoboard.

\textbf{Circuitos integrados utilizados:}
\begin{itemize}[nosep]
    \item CD4069 (hexa inversor) --- para los inversores del comparador de
    contraseña.
    \item CD4081 (AND doble de 2 entradas) / CD4073 (AND triple de 3 entradas)
    --- para las compuertas AND del comparador y de las señales $Trigger$ y
    $S_0$.
    \item 74HC74 (flip-flop D doble) --- para el módulo de cambio de estado.
\end{itemize}

\textbf{Incidencias resueltas durante el armado:}

\begin{enumerate}
    \item \emph{Compuertas quemadas durante el cableado de switches.} Se
    detectó un cableado diagonal que puenteaba columnas no relacionadas de la
    protoboard, generando una ruta de baja impedancia entre Vdd y GND. Se
    corrigió verificando el trayecto exacto de cada conexión y asegurando que
    cada entrada tuviera su resistencia de \emph{pull-down} (10\,k$\Omega$)
    hacia GND, evitando entradas flotantes.

    \item \emph{Ambos LEDs permanecían encendidos simultáneamente.} Se
    identificó que los pines \texttt{PRE} y \texttt{CLR} del 74HC74 son activos
    en bajo (a diferencia del CD4013, cuyos pines equivalentes \texttt{SET} y
    \texttt{RESET} son activos en alto). Al conectar \texttt{PRE} y \texttt{CLR}
    a GND —replicando la lógica del CD4013— se forzaba una condición inválida de
    preset y clear simultáneos, obligando ambas salidas ($Q$ y $Q'$) a estado
    alto al mismo tiempo. La corrección consistió en conectar ambos pines a Vcc
    para desactivarlos correctamente.
\end{enumerate}

Al finalizar el bloque, los tres módulos (comparador de contraseña, toggle de
estado de puerta y anuncio en display) fueron verificados de forma independiente
en la simulación, confirmando el comportamiento esperado según las ecuaciones
booleanas derivadas el 3 de septiembre.

\subsection*{Simulación funcional en Logisim Evolution}

Se simuló el módulo de comparación de contraseña en Logisim Evolution, verificando
el comportamiento esperado de las señales $Eq_1$, $Eq_2$, $Eq_3$, $Correcto$
(etiquetada $C$ en el esquema), $S_0$, $S_1$ y las salidas hacia los LEDs verde y
rojo del flip-flop D. La simulación confirmó que el circuito funciona
correctamente según las ecuaciones booleanas derivadas el 3 de septiembre. El
archivo de proyecto (\texttt{.circ}) se incluye en el \texttt{.zip} de la
entrega.

\begin{figure}[h]
    \centering
    \includegraphics[width=\textwidth]{logisim_comparador.png}
    \caption{Esquemático simulado en Logisim Evolution: comparador de
    contraseña ($Eq_1$--$Eq_3$ y $Correcto$), decodificador de anuncios
    ($S_0$) y disparo del flip-flop D hacia los LEDs de estado de la puerta.}
\end{figure}

\vspace{0.3cm}
\hrule
\vspace{0.3cm}

\section*{Pendientes para próximas sesiones}

\begin{itemize}[nosep]
    \item Armado físico en protoboard con la disciplina de niveles CMOS
    (0V--5V), usando la simulación de Logisim como referencia de comparación.
    \item Documentar el debounce del pulsador de confirmar como decisión de
    diseño pendiente de implementación.
    \item Documentar formalmente el criterio usado para descartar los estados
    de apertura/cierre en el display, para incluirlo en el paper.
\end{itemize}

\end{document}
